\documentclass{article}
\usepackage[UTF8]{ctex}
\usepackage{amsmath}
\usepackage{amssymb}
\usepackage{geometry}
\geometry{a4paper, margin=2.5cm}

\title{力控制误差收敛性证明：$K_{fp}\dot{F}_e + K_{fi}F_e = 0$}
\author{}
\date{}

\begin{document}

\maketitle

\section{问题陈述}

在力控制中，当存在非零但恒定的力干扰时（例如，由于重力模型$\tilde{g}(\theta)$不准确），误差动力学方程为：
\begin{equation}
K_{fp}\dot{F}_e + K_{fi}F_e = 0, \tag{67}
\end{equation}
其中：
\begin{itemize}
\item $F_e = F_d - F_{\text{tip}}$是力误差（期望力减去实际力）
\item $K_{fp}$是正定比例增益矩阵
\item $K_{fi}$是正定积分增益矩阵
\end{itemize}

我们需要证明：对于正定的$K_{fp}$和$K_{fi}$，$F_e$会收敛到零。

\section{背景：误差动力学方程的推导}

\subsection{完美重力建模的情况}

在完美重力建模的情况下，力控制器为：
\begin{equation}
\tau = \tilde{g}(\theta) + J^T(\theta)\left[F_d + K_{fp}F_e + K_{fi}\int_0^t F_e(t)dt\right],
\end{equation}
其中$\tilde{g}(\theta)$是重力补偿项。

代入动力学方程（忽略加速度和速度项）：
\begin{equation}
g(\theta) + J^T(\theta)F_{\text{tip}} = \tau,
\end{equation}
我们得到误差动力学：
\begin{equation}
K_{fp}F_e + K_{fi}\int_0^t F_e(t)dt = 0. \tag{11.52}
\end{equation}

\subsection{存在恒定干扰的情况}

当重力模型不准确时，存在恒定干扰$d$（常数向量），误差动力学变为：
\begin{equation}
K_{fp}F_e + K_{fi}\int_0^t F_e(t)dt = d, \tag{11.52'}
\end{equation}
其中$d$是恒定向量（例如，$d = g(\theta) - \tilde{g}(\theta)$在稳态时的值）。

\subsection{求导得到微分方程}

对方程(11.52')两边求时间导数：
\begin{equation}
\frac{d}{dt}\left[K_{fp}F_e + K_{fi}\int_0^t F_e(t)dt\right] = \frac{d}{dt}[d].
\end{equation}

由于$d$是常数，$\frac{d}{dt}[d] = 0$，且$\frac{d}{dt}\int_0^t F_e(t)dt = F_e$，我们得到：
\begin{equation}
K_{fp}\dot{F}_e + K_{fi}F_e = 0. \tag{67}
\end{equation}

\section{证明：$F_e$收敛到零}

\subsection{方法1：直接求解微分方程}

方程(67)是一个一阶线性齐次微分方程。为了求解，我们首先将其写成标准形式。

\subsubsection{标量情况（简化理解）}

首先考虑标量情况，其中$K_{fp}$和$K_{fi}$是正标量：
\begin{equation}
K_{fp}\dot{F}_e + K_{fi}F_e = 0.
\end{equation}

重写为：
\begin{equation}
\dot{F}_e = -\frac{K_{fi}}{K_{fp}}F_e.
\end{equation}

这是一个标准的一阶线性齐次微分方程，解为：
\begin{equation}
F_e(t) = F_e(0)e^{-\frac{K_{fi}}{K_{fp}}t}.
\end{equation}

\textbf{收敛性}：
\begin{itemize}
\item 由于$K_{fp} > 0$和$K_{fi} > 0$，我们有$\frac{K_{fi}}{K_{fp}} > 0$
\item 因此$e^{-\frac{K_{fi}}{K_{fp}}t} \to 0$当$t \to \infty$
\item 所以$F_e(t) \to 0$当$t \to \infty$
\end{itemize}

\subsubsection{向量情况（一般情况）}

对于向量情况，方程(67)是矩阵微分方程：
\begin{equation}
K_{fp}\dot{F}_e + K_{fi}F_e = 0,
\end{equation}
其中$K_{fp}$和$K_{fi}$是正定矩阵，$F_e \in \mathbb{R}^6$。

重写为：
\begin{equation}
\dot{F}_e = -K_{fp}^{-1}K_{fi}F_e. \tag{67'}
\end{equation}

\textbf{关键观察}：矩阵$-K_{fp}^{-1}K_{fi}$的特征值决定了系统的稳定性。

\subsection{方法2：使用Lyapunov稳定性理论}

我们构造一个Lyapunov函数来证明稳定性。

\subsubsection{构造Lyapunov函数}

定义Lyapunov函数：
\begin{equation}
V(F_e) = \frac{1}{2}F_e^T K_{fp}^{-1}K_{fi}F_e.
\end{equation}

\textbf{正定性}：
\begin{itemize}
\item 由于$K_{fp}$和$K_{fi}$都是正定矩阵，$K_{fp}^{-1}$也是正定的
\item 两个正定矩阵的乘积$K_{fp}^{-1}K_{fi}$的特征值都是正的（虽然矩阵本身可能不对称）
\item 更准确地说，$K_{fp}^{-1}K_{fi}$是正定的（可以通过相似变换证明）
\item 因此$V(F_e) \geq 0$，且$V(F_e) = 0$当且仅当$F_e = 0$
\end{itemize}

实际上，更简单的方法是使用：
\begin{equation}
V(F_e) = \frac{1}{2}F_e^T K_{fp}F_e.
\end{equation}

这个函数显然是正定的，因为$K_{fp}$是正定矩阵。

\subsubsection{计算$\dot{V}$}

对$V$求时间导数：
\begin{align}
\dot{V} &= \frac{d}{dt}\left[\frac{1}{2}F_e^T K_{fp}F_e\right] \\
        &= \frac{1}{2}\dot{F}_e^T K_{fp}F_e + \frac{1}{2}F_e^T K_{fp}\dot{F}_e \\
        &= F_e^T K_{fp}\dot{F}_e \quad \text{（因为$K_{fp}$是对称的）}
\end{align}

从方程(67)，我们有$K_{fp}\dot{F}_e = -K_{fi}F_e$，因此：
\begin{equation}
\dot{V} = F_e^T(-K_{fi}F_e) = -F_e^T K_{fi}F_e.
\end{equation}

\subsubsection{证明$\dot{V} \leq 0$}

由于$K_{fi}$是正定矩阵，对于任何非零向量$F_e$，我们有：
\begin{equation}
F_e^T K_{fi}F_e > 0.
\end{equation}

因此：
\begin{equation}
\dot{V} = -F_e^T K_{fi}F_e \leq 0,
\end{equation}
且$\dot{V} = 0$当且仅当$F_e = 0$。

\subsubsection{应用Lyapunov稳定性定理}

根据Lyapunov稳定性定理：
\begin{itemize}
\item $V(F_e) > 0$对所有$F_e \neq 0$成立，且$V(0) = 0$
\item $\dot{V}(F_e) < 0$对所有$F_e \neq 0$成立，且$\dot{V}(0) = 0$
\end{itemize}

因此，平衡点$F_e = 0$是\textbf{全局渐近稳定}的，即：
\begin{equation}
\lim_{t \to \infty} F_e(t) = 0,
\end{equation}
从任何初始条件$F_e(0)$开始。

\subsection{方法3：特征值分析}

从方程(67')：
\begin{equation}
\dot{F}_e = -K_{fp}^{-1}K_{fi}F_e = -A F_e,
\end{equation}
其中$A = K_{fp}^{-1}K_{fi}$。

\subsubsection{矩阵$A$的性质}

\textbf{关键性质}：虽然$A = K_{fp}^{-1}K_{fi}$本身可能不对称，但我们可以证明它的所有特征值都有正实部。

考虑广义特征值问题：
\begin{equation}
K_{fi}v = \lambda K_{fp}v,
\end{equation}
其中$\lambda$是广义特征值，$v$是广义特征向量。

这等价于：
\begin{equation}
K_{fp}^{-1}K_{fi}v = \lambda v,
\end{equation}
即$\lambda$是$A = K_{fp}^{-1}K_{fi}$的特征值。

\textbf{正定性}：
\begin{itemize}
\item 由于$K_{fp}$和$K_{fi}$都是正定矩阵
\item 广义特征值$\lambda$都是正的（这是正定矩阵对的标准性质）
\item 因此$A$的所有特征值都是正的
\end{itemize}

\subsubsection{系统矩阵的特征值}

系统矩阵是$-A = -K_{fp}^{-1}K_{fi}$，其特征值是$A$的特征值的负值。

由于$A$的所有特征值都是正的，$-A$的所有特征值都是负的。

\subsubsection{收敛性结论}

对于线性系统$\dot{F}_e = -A F_e$，如果$-A$的所有特征值都有负实部，则系统是稳定的，且：
\begin{equation}
F_e(t) = e^{-At}F_e(0) \to 0 \quad \text{当} \quad t \to \infty.
\end{equation}

更准确地说，如果$-A$的所有特征值$\lambda_i$都满足$\text{Re}(\lambda_i) < 0$，则：
\begin{equation}
\|F_e(t)\| \leq Ce^{-\alpha t}\|F_e(0)\|,
\end{equation}
其中$C > 0$是常数，$\alpha = \min_i \text{Re}(-\lambda_i) > 0$。

因此$F_e(t)$指数收敛到零。

\subsection{方法4：能量耗散观点}

从物理角度看，方程(67)描述了一个能量耗散过程。

\subsubsection{能量函数}

定义"误差能量"：
\begin{equation}
E(t) = \frac{1}{2}F_e^T(t)K_{fp}F_e(t).
\end{equation}

\subsubsection{能量变化率}

\begin{align}
\dot{E} &= F_e^T K_{fp}\dot{F}_e \\
        &= F_e^T(-K_{fi}F_e) \quad \text{（从方程(67)）} \\
        &= -F_e^T K_{fi}F_e \leq 0.
\end{align}

\subsubsection{能量耗散}

\begin{itemize}
\item 能量$E(t)$单调递减：$\dot{E} \leq 0$
\item 能量有下界：$E(t) \geq 0$
\item 因此能量收敛到某个值：$\lim_{t \to \infty} E(t) = E_{\infty} \geq 0$
\end{itemize}

\subsubsection{证明$E_{\infty} = 0$}

假设$E_{\infty} > 0$，则$F_e(t)$不收敛到零，即存在$\epsilon > 0$使得$\|F_e(t)\| \geq \epsilon$对所有$t$成立。

那么：
\begin{equation}
\dot{E} = -F_e^T K_{fi}F_e \leq -\lambda_{\min}(K_{fi})\|F_e\|^2 \leq -\lambda_{\min}(K_{fi})\epsilon^2 < 0,
\end{equation}
其中$\lambda_{\min}(K_{fi}) > 0$是$K_{fi}$的最小特征值。

这意味着能量以恒定速率减少，与$E(t)$有下界矛盾。因此$E_{\infty} = 0$，即$F_e(t) \to 0$。

\section{收敛速率}

\subsection{指数收敛}

从特征值分析，我们知道$F_e(t)$指数收敛到零：
\begin{equation}
\|F_e(t)\| \leq Ce^{-\alpha t}\|F_e(0)\|,
\end{equation}
其中收敛速率$\alpha$取决于$K_{fp}$和$K_{fi}$的选择。

\subsection{时间常数}

对于标量情况，时间常数为：
\begin{equation}
\tau = \frac{K_{fp}}{K_{fi}}.
\end{equation}

误差在时间$t = \tau$时衰减到初始值的$1/e \approx 37\%$。

\subsection{设计考虑}

\begin{itemize}
\item \textbf{快速收敛}：增大$K_{fi}$或减小$K_{fp}$可以加快收敛
\item \textbf{稳定性}：但过大的$K_{fi}$可能导致系统不稳定（在实际系统中，由于延迟、噪声等因素）
\item \textbf{平衡}：需要在收敛速度和稳定性之间找到平衡
\end{itemize}

\section{总结}

我们通过多种方法证明了：对于正定的$K_{fp}$和$K_{fi}$，误差动力学方程
\begin{equation}
K_{fp}\dot{F}_e + K_{fi}F_e = 0
\end{equation}
的解$F_e(t)$会\textbf{全局渐近收敛到零}，即：
\begin{equation}
\lim_{t \to \infty} F_e(t) = 0.
\end{equation}

\textbf{证明方法}：
\begin{enumerate}
\item \textbf{直接求解}：标量情况下直接求解，向量情况下使用矩阵指数
\item \textbf{Lyapunov方法}：构造Lyapunov函数$V = \frac{1}{2}F_e^T K_{fp}F_e$，证明$\dot{V} < 0$
\item \textbf{特征值分析}：证明系统矩阵的所有特征值都有负实部
\item \textbf{能量耗散}：从能量角度证明能量单调递减到零
\end{enumerate}

\textbf{关键条件}：
\begin{itemize}
\item $K_{fp}$必须是正定矩阵
\item $K_{fi}$必须是正定矩阵
\item 这两个条件保证了系统的全局渐近稳定性
\end{itemize}

\end{document}

